% Part: model-theory
% Chapter: models-of-arithmetic
% Section: models-of-pa

\documentclass[../../../include/open-logic-section]{subfiles}

\begin{document}

\olfileid{mod}{mar}{mpa}
\section{نماذج $\Th{PA}$}

\begin{explain}
كل نموذج غير قياسي لـ~$\Th{TA}$ هو أيضًا نموذج لـ~$\Th{PA}$. ونعلم أن
لـ~$\Th{TA}$ نماذج غير قياسية، ومن ثم أن لـ~$\Th{PA}$ نماذج كذلك. ونعلم
أيضًا أن هذه النماذج غير القياسية تحتوي «أعدادًا» غير قياسية، أي
!!{element}s من المجال تقع «بعد» جميع «الأعداد» القياسية. ولكن كيف
تنتظم؟ وكم عددها؟ لقد رأينا أن نماذج النظرية الأضعف~$\Th{Q}$ قد تحتوي
عددًا غير قياسي واحدًا لا غير. غير أن هذه الـ!!{structure}s البسيطة ليست
نماذج لـ$\Th{PA}$ ولا لـ$\Th{TA}$.

ومفتاح فهم بنية نماذج $\Th{PA}$ أو $\Th{TA}$ هو النظر في الوقائع التي
تكون !!{derivable} في هاتين النظريتين. فمثلًا تثبت $\Th{PA}$ نفسها
$\lforall[x][\eq/[x][x']]$ و
$\lforall[x][\lforall[y][\eq[(x+y)][(y+x)]]]$، وهذا يستبعد البنى البسيطة
(التي تكون فيها هذه الـ!!{sentence}s كاذبة) من أن تكون نماذج لـ~$\Th{PA}$.

افترض أن~$\Struct{M}$ نموذج لـ~$\Th{PA}$. فإذا كان
$\Th{PA} \Proves !A$، فإن $\Sat{M}{!A}$. ولنستعمل مرة أخرى $\nszero$
اختصارًا لـ$\Assign{\Obj{0}}{M}$، و$\nssucc$ لـ$\Assign{\prime}{M}$،
و$\nsplus$ لـ$\Assign{+}{M}$، و$\nstimes$ لـ$\Assign{\times}{M}$،
و$\nsless$ لـ$\Assign{<}{M}$. وعندئذ تقرر كل !!{sentence}~$!A$ شرطًا ما
على $\nszero$ و$\nssucc$ و$\nsplus$ و$\nstimes$ و$\nsless$، وإذا كان
$\Sat{M}{!A}$ فلا بد أن يتحقق ذلك الشرط. فمثلًا، إذا كان
$\Sat{M}{!Q_1}$، أي
$\Sat{M}{\lforall[x][\lforall[y][(\eq[x'][y'] \lif \eq[x][y])]]}$،
فلا بد أن تكون $\nssucc$~!!{injective}.
\end{explain}

\begin{prop}
تكون $\nsless$ في $\Struct{M}$ ترتيبًا خطيًا صارمًا، أي إنها تحقق:
\begin{enumerate}
\item ليس $x \nsless x$ لأي~$x \in \Domain{M}$.
\item إذا كان $x \nsless y$ و$y \nsless z$، فإن $x \nsless z$.
\item لأي $x \neq y$، إما $x \nsless y$ وإما $y \nsless x$.
\end{enumerate}
\end{prop}

\begin{proof}
تثبت $\Th{PA}$ ما يأتي:
\begin{enumerate}
\item $\lforall[x][\lnot x < x]$
\item $\lforall[x][\lforall[y][\lforall[z][((x < y \land y < z) \lif x < z)]]]$
\item $\lforall[x][\lforall[y][((x < y \lor y < x) \lor \eq[x][y])]]$
\end{enumerate}
\end{proof}

\begin{prop}
\ollabel{prop:M-discrete} $\nszero$ أصغر !!{element} في~$\Domain{M}$
بحسب ترتيب~$\nsless$. ولكل~$x$، يكون $x \nsless x^\nssucc$، ويكون
$x^\nssucc$ أصغر !!{element} بحسب~$\nsless$ له هذه الخاصية. ولكل
$x \neq \nszero$ يوجد $y$ وحيد بحيث $y^\nssucc = x$. (نسمي $y$
«سلف»~$x$ في~$\Struct{M}$، ونرمز إليه بـ~$^\nssucc x$.)
\end{prop}

\begin{proof}
تمرين.
\end{proof}

\begin{prob}
جد !!{sentence}s في~$\Lang{L_A}$ تكون !!{derivable} في~$\Th{PA}$
(ومن ثم صادقة في~$\Struct{N}$) وتضمن خصائص $\nszero$ و$\nssucc$
و$\nsless$ الواردة في
\olref[mod][mar][mpa]{prop:M-discrete}.
\end{prob}

\begin{prop}
كل الـ!!{element}s القياسية في~$\Struct{M}$ أصغر (بحسب~$\nsless$) من
جميع الـ!!{element}s غير القياسية.
\end{prop}

\begin{proof}
سنستعمل $n$ اختصارًا لـ$\Value{\num{n}}{M}$، وهو !!{element} قياسي
في~$\Struct{M}$. تثبت $\Th{Q}$ نفسها، لكل~$n \in \Nat$، أن
$\lforall[x][(x < \num{n}' \lif (\eq[x][\num{0}] \lor
  \eq[x][\num{1}] \lor \dots \lor \eq[x][\num{n}]))]$. ولا توجد
!!{element}s أصغر من $\nszero$ بحسب~$\nsless$. لذلك إذا كان $n$
قياسيًا وكان $x$ غير قياسي، استحال أن يكون $x \nsless n$. وبالتعريف،
العنصر غير القياسي هو ما لا يساوي $\Value{\num{n}}{M}$ لأي
$n \in \Nat$، ومن ثم يكون $x \neq n$ أيضًا. ولما كانت $\nsless$
ترتيبًا خطيًا، فلا بد أن يكون $n \nsless x$.
\end{proof}

\begin{prop}
كل !!{element} غير قياسي~$x$ في~$\Domain{M}$ هو عنصر في المجموعة الجزئية
\[
\dots ^{\nssucc\nssucc\nssucc}x \nsless ^{\nssucc\nssucc}x \nsless
^{\nssucc}x \nsless x \nsless x^{\nssucc} \nsless x^{\nssucc\nssucc}
\nsless x^{\nssucc\nssucc\nssucc} \nsless \dots
\]
ونسمي هذه المجموعة الجزئية \emph{كتلة~$x$}، ونرمز إليها بـ$[x]$.
وليس لها أصغر عنصر ولا أكبر !!{element}. ويمكن وصفها بأنها مجموعة
كل $y \in \Domain{M}$ بحيث يوجد عدد قياسي~$n$ يحقق
$x \nsplus n = y$ أو $y \nsplus n = x$.
\end{prop}

\begin{proof}
من الواضح أن مثل هذه المجموعة~$[x]$ موجود دائمًا، لأن كل !!{element}
غير قياسي~$y$ في~$\Domain{M}$ له خلف وحيد~$y^\nssucc$ وسلف وحيد
$^\nssucc y$. ولأي !!{element}s متعاقبة $y$ و$y^\nssucc$ يكون
$y \nsless y^\nssucc$، ويكون $y^\nssucc$ أصغر !!{element} في
$\Domain{M}$ بحسب~$\nsless$ يكون $y$ أصغر منه بحسب~$\nsless$.
ولما كان دائمًا $^\nssucc y \nsless y$ و$y \nsless y^\nssucc$، لم
يكن لـ$[x]$ أصغر عنصر ولا أكبر !!{element}. وإذا كان
$y \in [x]$، فإن $x \in [y]$، إذ يكون حينئذ إما
$y^{\nssucc\dots\nssucc} = x$ وإما
$x^{\nssucc\dots\nssucc} = y$. وإذا كان
$y^{\nssucc\dots\nssucc} = x$ (وفيه $n$ من رموز~$\nssucc$)، فإن
$y \nsplus n = x$، والعكس صحيح، لأن
$\Th{PA} \Proves \lforall[x][\eq[x^{\prime\dots\prime}][(x +
    \num{n})]]$ (إذا كان $n$ عدد رموز~$\prime$).
\end{proof}

\begin{prop}
إذا كان $[x] \neq [y]$ و$x \nsless y$، فإن $u \nsless v$ لكل
$u \in [x]$ ولكل $v \in [y]$.
\end{prop}

\begin{proof}
لاحظ أن
$\Th{PA} \Proves \lforall[x][\lforall[y][(x < y \lif (x' < y
    \lor x' = y))]]$. ولذلك إذا كان $u \nsless v$، كان أيضًا
$u \nsplus n^\nssucc \nsless v$ لأي~$n$ متى كان $[u] \neq [v]$.

وكل $u \in [x]$ أصغر من~$y$ بحسب~$\nsless$: فبحسب الفرض
$x \nsless y$. فإذا كان $u \nsless x$، كان $u \nsless y$ بالتعدي.
وإذا كان $x \nsless u$ مع $u \in [x]$، كان
$u = x\nsplus n^\nssucc$ لبعض~$n$، ومن ثم $u \nsless y$ بالواقعة التي
أثبتناها للتو.

افترض الآن أن $v \in [y]$ أصغر من~$y$ بحسب~$\nsless$، أي إن
$v \nsplus m^\nssucc = y$ لبعض $m$ قياسي. وهذا يستبعد
$v \nsless x$، وإلا لكان $y = v \nsplus m^\nssucc \nsless x$.
ومن الواضح أيضًا أن $x \neq v$، وإلا لكان
$x \nsplus m^\nssucc = v \nsplus m^\nssucc = y$، ولحصلنا على
$[x] = [y]$. إذن $x \nsless v$. ولكن عندئذ يكون أيضًا
$x \nsplus n^\nssucc \nsless v$ لأي~$n$. ومن ثم، إذا كان
$x \nsless u$ و$u \in [x]$، كان $u \nsless v$. وإذا كان
$u \nsless x$، كان $u \nsless v$ بالتعدي.

وأخيرًا، إذا كان $y \nsless v$، فإن $u \nsless v$ لأننا بينا أن
$u \nsless y$، ولدينا $y \nsless v$.
\end{proof}

\begin{cor}
إذا كان $[x] \neq [y]$، فإن $[x] \cap [y] = \emptyset$.
\end{cor}

\begin{proof}
افترض أن $z \in [x]$ وأن $x \nsless y$. عندئذ يكون $z \nsless u$
لكل $u \in [y]$. ولو كان $z \in [y]$، لحصلنا على $z \nsless z$.
والحجة مماثلة إذا كان $y \nsless x$.
\end{proof}

\begin{explain}
يعني هذا أن الكتل نفسها يمكن ترتيبها على نحو يحترم~$\nsless$:
$[x] \nsless [y]$ إذا وفقط إذا كان $x \nsless y$، أو، على نحو مكافئ،
إذا كان $u \nsless v$ لأي $u \in [x]$ و$v \in [y]$. ومن الواضح أن
الكتلة القياسية~$[0]$ هي أصغر الكتل. وهي لا تتقاطع مع أي كتلة غير
قياسية، ولا تتقاطع أي كتلتين غير قياسيتين أيضًا. وعلى وجه التحديد، لا
يمكن «بلوغ» كتلة مختلفة بأخذ الأخلاف أو الأسلاف المتكررة.
\end{explain}

\begin{prop}
إذا كان $x$ و$y$ غير قياسيين، فإن $x \nsless x \nsplus y$ و
$x \nsplus y \notin [x]$.
\end{prop}

\begin{proof}
إذا كان $y$ غير قياسي، فإن $y \neq \nszero$. ولدينا
$\Th{PA} \Proves \lforall[x][(y \neq \Obj{0} \lif x < (x+y))]$. افترض الآن أن
$x \nsplus y \in [x]$. ولما كان $x \nsless x \nsplus y$، فلا بد أن
يكون $x \nsplus n^\nssucc = x \nsplus y$. لكن لدينا
$\Th{PA} \Proves \lforall[x][\lforall[y][\lforall[z][(\eq[(x+y)][(x+z)] \lif y = z)]]]$
(قانون الاختزال للجمع). ويعني هذا أن $y = n^\nssucc$ لبعض $n$ قياسي؛
لكن المفروض أن $y$ غير قياسي.
\end{proof}

\begin{prop}
لا توجد كتلة غير قياسية صغرى.
\end{prop}

\begin{proof}
لدينا $\Th{PA} \Proves \lforall[x][\lexists[y][(\eq[(y+y)][x] \lor
      \eq[(y+y)'][x])]]$، أي إن كل $x$ يقبل القسمة على~$2$ (وربما
مع باق مقداره~$1$). وإذا كان $x$ غير قياسي، كان $y$ غير قياسي أيضًا.
وبحسب القضية السابقة، يكون $y \nsless y \nsplus y$ و
$y \nsplus y \notin [y]$. وكذلك يكون
$y \nsless (y \nsplus y)^\nssucc$ و
$(y \nsplus y)^\nssucc \notin [y]$. لكن
$x = y \nsplus y$ أو $x = (y \nsplus y)^\nssucc$، ومن ثم
$y \nsless x$ و$y \notin [x]$.
\end{proof}

\begin{prop}
لا توجد كتلة كبرى.
\end{prop}

\begin{proof}
تمرين.
\end{proof}

\begin{prob}
بيّن أنه لا توجد كتلة كبرى في أي نموذج غير قياسي لـ~$\Th{PA}$.
\end{prob}

\begin{prop}
\ollabel{prop:blocks-dense}
ترتيب الكتل كثيف. أي إنه إذا كان $x \nsless y$ و$[x] \neq [y]$،
فثمة كتلة~$[z]$ تختلف عن كلتيهما وتقع بينهما.
\end{prop}

\begin{proof}
افترض أن $x \nsless y$. وكما سبق، يقبل $x \nsplus y$ القسمة على
اثنين (وربما مع باق): يوجد $z \in \Domain{M}$ بحيث إما
$x \nsplus y = z \nsplus z$ أو
$x \nsplus y = (z \nsplus z)^\nssucc$. والعنصر~$z$ هو
«متوسط» $x$ و$y$، ويكون $x \nsless z$ و$z \nsless y$.
\end{proof}

\begin{prob}
اكتب برهانًا مفصلًا لـ
\olref[mod][mar][mpa]{prop:blocks-dense}. أي !!{sentence} يجب أن
!!{derive}ها $\Th{PA}$ كي تضمن وجود~$z$؟ ولماذا يكون $x \nsless z$
و$z \nsless y$، ولماذا يكون $[x] \neq [z]$ و$[z] \neq [y]$؟
\end{prob}

\begin{explain}
وهكذا تكون الكتل غير القياسية مرتبة مثل الأعداد النسبية: فهي تشكل
ترتيبًا خطيًا كثيفًا !!a{denumerable} بلا طرفين. ويمكن إثبات أن كل ترتيبين
!!{denumerable} من هذا القبيل متشاكلان. وينتج من ذلك أنه لأي نموذجين
غير قياسيين !!{enumerable}~$\Struct{M}_1$ و$\Struct{M}_2$ للحساب الصادق،
يكون مختزلاهما إلى اللغة التي لا تحتوي إلا $<$ و$=$ متشاكلين. والواقع
أنه يمكن تعريف تشاكل~$h$ كما يأتي: الجزآن القياسيان من $\Struct{M}_1$
و$\Struct{M}_2$ متشاكلان مع النموذج القياسي~$\Struct{N}$، ومن ثم
أحدهما مع الآخر. والكتل التي يتألف منها الجزء غير القياسي مرتبة هي نفسها
مثل الأعداد النسبية، ومن ثم فهي متشاكلة؛ ويمكن مد تشاكل الكتل إلى تشاكل
\emph{داخل} الكتل بمقابلة عناصر اعتباطية في كل منها، ثم جعل صورة خلف~$x$
في~$\Struct{M}_1$ خلف صورة~$x$ في~$\Struct{M}_2$. لاحظ أنه \emph{لا} ينتج من
ذلك أن $\mathfrak{M}_1$ و$\mathfrak{M}_2$ متشاكلان في لغة الحساب الكاملة
(فالتشاكل نسبي دائمًا إلى !!a{language})، إذ توجد طرائق غير متشاكلة لتعريف
الجمع والضرب على $\Domain{M_1}$ و$\Domain{M_2}$. (وينتج هذا أيضًا من
مبرهنة شهيرة لفوت تنص على أن عدد النماذج القابلة للعد لنظرية تامة لا يمكن
أن يكون~2.)
\end{explain}
\end{document}
